%!TEX root = ../thesis.tex
\chapter{Related Work}
\label{chap:relatedwork}

\noindent
This chapter reviews three relevant lines of work: self-supervised pretraining applied to driving, cross-city and geographic generalization, and diffusion-based trajectory planning. The technical material on NAVSIM, the three SSL objectives, and the four planners used in the experiments is in Chapter~\ref{chap:background} and is not restated here.

% ============================================================================
\section{Self-Supervised Learning for Driving}
\label{sec:rw-ssl-driving}

Self-supervised pretraining has been applied extensively to driving perception and only sparsely to end-to-end planning. On the LiDAR and BEV side, PointContrast~\cite{xie-pointcontrast-2020}, Occupancy-MAE~\cite{min-occupancymae-2023}, BEV-MAE~\cite{lin-bevmae-2024}, UniPAD~\cite{yang-unipad-2024}, and AD-L-JEPA~\cite{zhu-adljepa-2025} show label-efficiency and detection gains with self-supervised pretext tasks, but report perception metrics rather than closed-loop planning scores. On the camera side, the closest precedents are ViDAR~\cite{yang-vidar-2024}, which improves UniAD's downstream tasks through an image-side pretext; GAIA-1~\cite{hu-gaia1-2023}, a large-scale video-and-action generative model that trades compute for breadth; Juneja et al.~\cite{juneja-dino-e2e-2024}, who show that a DINO-initialized backbone beats ImageNet pretraining on CARLA; and LAW~\cite{li-law-2024}, whose latent world-model loss is itself a self-supervised auxiliary on top of a supervised Swin-T encoder.

The most relevant recent result is Drive-JEPA~\cite{wang-drive-jepa}, which uses a V-JEPA video encoder together with multi-modal trajectory distillation and reaches strong pooled NAVSIM PDMS. Drive-JEPA pretraining mixes driving video from Japan, the United States, Singapore, and China, so its evaluation blends pretraining geography with downstream evaluation geography. That makes it informative about the in-distribution ceiling of SSL-initialized planning, and simultaneously motivates this thesis's focus on city-disjoint evaluation: pooled NAVSIM PDMS is no longer the right discriminator for a representation-centred argument. Plan-MAE~\cite{zhao-planmae-2025} is a related masked-autoencoding approach applied to joint prediction and planning.

The limitation of this literature, for the thesis's purpose, is twofold. Published SSL-for-driving results routinely mix pretraining geography with downstream evaluation geography, so transfer gains cannot be attributed cleanly to the SSL objective rather than to pretraining-data overlap. When multiple SSL objectives are compared, the comparison is typically confined to one downstream task and one backbone. The thesis addresses both by fixing the planner, varying only the backbone, and evaluating under a strict city-disjoint protocol.

% ============================================================================
\section{Cross-City and Geographic Generalization}
\label{sec:rw-dg}

Classical domain-generalization methods including IRM~\cite{arjovsky-irm-2019}, GroupDRO~\cite{sagawa-groupdro-2020}, and DANN~\cite{ganin-dann-2015} are reviewed extensively in~\cite{wang-domain-gen-survey-2022,zhou-domain-gen-survey-2023}. DomainBed~\cite{gulrajani-domainbed-2021} reported that empirical risk minimization matches or exceeds about twenty DG methods on seven benchmarks under careful model selection, and this result is why the thesis changes backbone initialization rather than the DG training procedure.

Most driving-specific DG work concentrates on appearance-level shift such as weather, lighting, and season. ACDC~\cite{sakaridis-acdc-2021}, DAFormer~\cite{hoyer-daformer-2022}, HRDA~\cite{hoyer-hrda-2022}, DACS~\cite{tranheden-dacs-2021}, CyCADA~\cite{pan-cycada-2018}, and the CARLA train-town/test-town protocol~\cite{dosovitskiy-carla-2017} are representative. Geographic shift is studied earlier at the perception level: Chen et al.~\cite{chen-crosscity-2017} provide the first explicit cross-city UDA benchmark; Beery et al.~\cite{beery-cameratrap-2018} and de Vries et al.~\cite{devries-geobias-2019} document geographic accuracy drops in visual recognition; GeoNet~\cite{kalluri-geonet-2023} formalizes unsupervised adaptation across geographies.

Cross-city transfer at the planning level is much sparser. VLP~\cite{pan-vlp-2024} reports Boston-to-Singapore and Singapore-to-Boston planning L2 and collision rates, but without analysing directional asymmetry. Trajectory-prediction systems such as GTG~\cite{gtg-2025} pursue city-invariant representations at the prediction layer rather than at the full planning layer. Naeinian et al.~\cite{naeinian-cross-city-2026} is the direct precursor to this thesis: it introduces strict city-disjoint evaluation across the four NAVSIM cities for both open-loop nuScenes and closed-loop NAVSIM, and it documents a directional asymmetry in which Boston-to-Singapore transfer is materially harder than Singapore-to-Boston. This thesis extends that protocol, adds a NAVSIM v1.1 re-evaluation, broadens the backbone comparison, integrates DiffusionDrive, and adds the mechanistic analysis in Chapter~\ref{chap:experiments}.

% ============================================================================
\section{Diffusion-Based Trajectory Planning}
\label{sec:rw-diffusion}

Diffusion models~\cite{ho-ddpm-2020} reverse a Gaussian noising process and can model multi-modal action distributions without the mode averaging that regression heads produce. DDIM sampling~\cite{song-ddim-2021}, classifier-free guidance~\cite{ho-cfg-2022}, and consistency models~\cite{song-consistency-2023} reduce the number of denoising steps enough to make diffusion tractable for real-time planning.

In robotics, Diffuser~\cite{janner-diffuser-2022} first denoised entire state-action trajectories with classifier guidance for rewards. Diffusion Policy~\cite{chi-diffusion-policy-2023} applied conditional DDPMs to visuomotor manipulation and reported a $46.9\%$ relative improvement over prior behavioural cloning across twelve tasks, which is the clearest published evidence that diffusion heads preserve multi-modal structure that regression heads collapse. MotionDiffuser~\cite{jiang-motiondiffuser-2023} extends diffusion to multi-agent motion prediction, and anchor-based alternatives such as MTR~\cite{shi-mtr-2022} and MultiPath++~\cite{varadarajan-multipathpp-2022} introduce learned latent anchors that inform the DiffusionDrive head used here.

For planning on NAVSIM, DiffusionDrive~\cite{jia-diffusiondrive-2025} is the reference diffusion planner and the one integrated with SSL backbones in Chapter~\ref{chap:experiments}. Related generative NAVSIM planners include GoalFlow~\cite{xing-goalflow-2025} and TransDiffuser~\cite{jiang-transdiffuser-2025}, which report high pooled PDMS but do not evaluate under city-disjoint splits. This thesis uses DiffusionDrive as the generative planner in both pooled mixed-city and city-disjoint evaluations. The contribution is bounded: prior NAVSIM work reviewed here does not isolate planner-head variation under a controlled cross-city holdout in the same way.
